\chapter{Distillation for Efficient Reasoning}\label{ch:distill}

Reinforcement learning taught a model to reason by trial, error and a verifier
(\cref{ch:grpo,ch:grpo2}), and \cref{ch:rlsys} showed how expensive that is: every optimizer
step pays for thousands of long generations. \emph{Distillation} takes a cheaper route. A
strong \emph{teacher} reasoning model writes out worked solutions once; a smaller
\emph{student} is then trained, with plain supervised learning, to reproduce them. The
DeepSeek-R1 report popularised this recipe for reasoning: the large RL-trained model produced
traces from which a family of much smaller dense models was fine-tuned, and the authors
reported that these distilled students outperformed comparable small models trained with RL
alone.

From a systems point of view distillation is attractive for a simple reason: it moves almost
all of the expensive inference \emph{out of the training loop}. Teacher generation becomes an
offline, embarrassingly parallel batch job whose output is a file; student training is ordinary
fine-tuning on that file, with fixed data, reproducible epochs and none of RL's
generator/trainer choreography. The source book reports that its full distillation run of
Qwen3-0.6B fits in roughly 15\,GB of memory and finishes in about three hours on a desktop
GB10-class machine, whereas a few GRPO rounds on the same hardware took far longer and several
times the memory. This chapter builds that pipeline end to end and then asks the systems
questions it raises: how to store soft targets, how to train on long traces, and how to make
reasoning \emph{cheaper} rather than merely smaller.

\begin{figure}[t]
\centering
\resizebox{\linewidth}{!}{%
\begin{tikzpicture}[
  big/.style={draw=sysblue,fill=sysbluebg,rounded corners=3pt,minimum width=2.9cm,
    minimum height=1.1cm,align=center,font=\small\bfseries},
  sub/.style={draw=sysblue!60,fill=white,rounded corners=2pt,minimum width=2.6cm,
    minimum height=0.65cm,align=center,font=\scriptsize},
  num/.style={circle,fill=sysblue,text=white,font=\bfseries\scriptsize,inner sep=1.3pt},
  arr/.style={-{Stealth[length=2.3mm]},thick,sysblue!70}]
\node[big] (s1) at (0,0) {Why distil?};
\node[big] (s2) at (4.0,0) {Build the\\dataset};
\node[big] (s3) at (8.0,0) {Train the\\student};
\node[big] (s4) at (12.0,0) {Evaluate \&\\look ahead};
\foreach \n/\k in {s1/1,s2/2,s3/3,s4/4} \node[num] at (\n.north west) {\k};
\draw[arr] (s1)--(s2); \draw[arr] (s2)--(s3); \draw[arr] (s3)--(s4);
\node[sub] (a1) at (0,-1.3) {hard vs.\ soft};
\node[sub] (a2) at (0,-2.1) {teacher / student};
\node[sub] (b1) at (4.0,-1.3) {2a generate traces};
\node[sub] (b2) at (4.0,-2.1) {2b load records};
\node[sub] (b3) at (4.0,-2.9) {2c tokenise, filter, split};
\node[sub] (c1) at (8.0,-1.3) {3a load base model};
\node[sub] (c2) at (8.0,-2.1) {3b answer-only loss};
\node[sub] (c3) at (8.0,-2.9) {3c training loop};
\node[sub] (d1) at (12.0,-1.3) {MATH-500 accuracy};
\node[sub] (d2) at (12.0,-2.1) {soft KD, long seqs};
\node[sub] (d3) at (12.0,-2.9) {efficient reasoning};
\end{tikzpicture}}
\caption{Roadmap. Stages 1--4 follow the distillation workflow (motivation, dataset
construction, student training, evaluation); the lower boxes are the concrete steps of each
stage. Stage 4 also contains this book's systems extensions.}
\label{fig:distill:roadmap}
\end{figure}

% =====================================================================================
\section{What distillation is}\label{sec:distill:intro}\covers{RM 8.1}

Distillation trains a \emph{student} model to imitate a \emph{teacher}. For reasoning models
the thing being imitated is not only the final answer but the whole trace---the exploratory,
self-correcting chain of thought that precedes it. There are two families of methods,
distinguished by \emph{what} the student sees of the teacher (\cref{fig:distill:hardsoft}).

\paragraph{Hard distillation.} The teacher generates text; the student is trained with the
ordinary next-token cross-entropy loss on that text. The targets are the teacher's
\emph{sampled tokens}, i.e.\ one-hot vectors. Hard distillation is just supervised fine-tuning
(\cref{ch:instruct}) on synthetic data, and it needs nothing but the teacher's output text---so
it works with closed models behind an API and across model families with different
tokenizers. DeepSeek-R1's small distilled models were produced this way.

\paragraph{Soft distillation.} The student is trained to match the teacher's full
\emph{distribution} over the vocabulary at every position, by minimising a divergence such as
\begin{equation}
\mathcal{L}_{\text{KD}} = \frac{1}{|\mathcal{A}|}\sum_{t\in\mathcal{A}}
\KL\!\big(p_T(\cdot\mid y_{<t})\,\big\|\,p_S(\cdot\mid y_{<t})\big),
\label{eq:distill:kd}
\end{equation}
where $\mathcal{A}$ is the set of answer positions. A soft target says not only ``the next
token is \emph{Subtract}'' but also ``\emph{Multiply} was plausible and \emph{banana} was not''.
That extra information---Hinton, Vinyals and Dean (2015) called it \emph{dark
knowledge}---usually makes each training token more informative. The price is access to the
teacher's logits and a shared vocabulary.

\paragraph{Combined.} The classic knowledge-distillation objective mixes the two:
$\mathcal{L} = \alpha\,\mathcal{L}_{\text{CE}} + (1-\alpha)\,\tau^2\,\mathcal{L}_{\text{KD}}^{(\tau)}$,
where $\mathcal{L}_{\text{KD}}^{(\tau)}$ compares distributions softened by a temperature $\tau$
and the factor $\tau^2$ keeps gradient magnitudes comparable across temperatures.

\begin{figure}[t]
\centering
\resizebox{\linewidth}{!}{%
\begin{tikzpicture}[box/.style={draw,rounded corners,align=center,font=\small,minimum height=1cm},
  arr/.style={-{Stealth},thick}]
\node[box,fill=fmemory!25,minimum width=2.6cm] (T) at (0,0) {Teacher\\(large reasoner)};
\node[draw,dashed,rounded corners,font=\scriptsize,align=left,text width=3.3cm] (p) at (-4.2,0)
  {problem: ``A train leaves at 3pm \ldots{} when does it arrive?''};
\draw[arr] (p) -- (T);
% hard
\node[box,fill=white,text width=4.3cm,font=\scriptsize] (txt) at (5.2,1.5)
  {\texttt{<think>} Distance is 240\,km, speed 80\,km/h, so 3 h \ldots{} \texttt{</think>} \ldots{} $\boxed{6\text{pm}}$};
\node[font=\scriptsize\bfseries,anchor=south] at (txt.north) {hard: sampled text (token ids)};
\draw[arr] (T.east) -- ++(0.5,0) |- (txt.west);
% soft
\begin{scope}[shift={(3.2,-2.4)}]
\node[font=\scriptsize\bfseries,anchor=south] at (2.0,1.25) {soft: a distribution at every position};
\foreach \x/\h/\c in {0/0.35/gray,0.35/1.05/fcompute,0.7/0.6/fcompute,1.05/0.15/gray,1.4/0.1/gray,
  1.75/0.45/fcompute,2.1/0.08/gray,2.45/0.05/gray,2.8/0.2/gray,3.15/0.06/gray,3.5/0.03/gray}
  \fill[\c!60] (\x,0) rectangle (\x+0.28,\h);
\draw (-0.1,0) -- (3.9,0);
\node[font=\tiny,anchor=north] at (1.9,-0.02) {vocabulary ($V\approx152$k for Qwen3)};
\end{scope}
\draw[arr] (T.east) -- ++(0.5,0) |- (3.1,-2.0);
% student
\node[box,fill=fcompute!25,minimum width=2.6cm] (S) at (10.8,0) {Student\\(small model)};
\draw[arr] (txt.east) -| node[pos=0.25,above,font=\scriptsize]{CE on tokens} (S.north);
\draw[arr] (7.3,-1.8) -| node[pos=0.25,below,font=\scriptsize]{KL to distribution} (S.south);
\end{tikzpicture}}
\caption{Hard distillation trains the student on the teacher's sampled tokens (one-hot
targets); soft distillation trains it to match the teacher's full next-token distribution,
which requires the teacher's logits and a shared vocabulary.}
\label{fig:distill:hardsoft}
\end{figure}

For LLMs hard distillation dominates in practice, for three reasons: frontier teachers rarely
expose full distributions; student and teacher must share a tokenizer for
\cref{eq:distill:kd} to even be defined; and soft targets are enormous (see the napkin math
below). This chapter therefore follows the DeepSeek-R1 recipe---hard distillation---for the main
pipeline, and returns to soft distillation as a systems problem in \cref{sec:distill:soft}.

\begin{napkinmath}[Text versus distributions]
A reasoning trace of 3{,}000 tokens stored as token ids is $3000\times 4$\,B $=12$\,KB. The
same trace with the full Qwen3 next-token distribution at every position in bf16 is
$3000 \times 151{,}936 \times 2$\,B $\approx 0.91$\,GB. For 12{,}000 training problems that is
$\approx$144\,MB of text versus $\approx$10.9\,TB of soft targets---a factor of about 76{,}000.
Keeping only the top 32 entries per position (fp16 value + int32 index) shrinks it to
$\approx$6.9\,GB, which is why practical soft distillation stores sparse top-$k$ targets or
runs the teacher online.
\end{napkinmath}

\begin{pitfall}[Terms of service]
Many commercial model providers' terms restrict using their outputs to train competing
models. Check the licence of every teacher before generating a distillation set. Open-weight
teachers (DeepSeek-R1, the larger Qwen3 models) are a common choice partly for this reason.
\end{pitfall}

% =====================================================================================
\section{Generating a distillation dataset}\label{sec:distill:generate}\covers{RM 8.2}

Our student is the Qwen3-0.6B \emph{base} model used throughout Part~VI, and our training
problems are the $\approx$12{,}000 problems of the MATH training split that do not overlap with
the MATH-500 test set---the same pool used for RLVR in \cref{ch:grpo,ch:grpo2}. What changes
is the source of the supervision (\cref{fig:distill:rlvr}). In RLVR the student samples its
own answers and a verifier compares each final answer with the reference to produce a scalar
reward. In distillation a teacher writes a complete solution and the student is trained to
reproduce it \emph{token by token}: the supervision is dense (every token is a target) and
does not require the student to ever find a correct answer by itself.

\begin{figure}[t]
\centering
\resizebox{\linewidth}{!}{%
\begin{tikzpicture}[box/.style={draw,rounded corners,align=center,font=\small,minimum height=0.9cm},
  arr/.style={-{Stealth},thick}]
\node[font=\small\bfseries,anchor=west] at (-1.2,1.3) {RLVR (Chapters 27--28)};
\node[box,fill=gray!10] (q1) at (0,0) {MATH train\\problem};
\node[box,fill=fcompute!25] (st1) at (3.6,0) {student\\samples $G$ answers};
\node[box,fill=fdata!25] (v) at (7.6,0) {verifier vs.\\reference answer};
\node[box,fill=white] (r) at (11.2,0) {scalar reward\\per answer};
\draw[arr] (q1)--(st1); \draw[arr] (st1)--(v); \draw[arr] (v)--(r);
\draw[arr,dashed] (r.south) -- ++(0,-0.45) -| node[pos=0.25,below,font=\scriptsize]{policy gradient (sparse signal)} (st1.south);
\node[font=\small\bfseries,anchor=west] at (-1.2,-2.0) {Distillation (this chapter)};
\node[box,fill=gray!10] (q2) at (0,-3.2) {MATH train\\problem};
\node[box,fill=fmemory!25] (te) at (3.6,-3.2) {teacher writes\\trace + answer};
\node[box,fill=white] (f) at (7.6,-3.2) {stored target\\sequence};
\node[box,fill=fcompute!25] (st2) at (11.2,-3.2) {student: CE on\\every target token};
\draw[arr] (q2)--(te); \draw[arr] (te)--node[above,font=\scriptsize]{offline}(f); \draw[arr] (f)--(st2);
\end{tikzpicture}}
\caption{Where the training signal comes from. RLVR compares the student's own final answer
with a reference; distillation compares the student's token predictions with a
teacher-written solution. The teacher runs once, offline.}
\label{fig:distill:rlvr}
\end{figure}

A practical virtue is visible in the figure: teacher generation is decoupled from training.
The source book generated its 12{,}000 traces by calling the 671B-parameter DeepSeek-R1 through
a hosted API (at a cost it reports as roughly \$50), and publishes the result so that readers
can skip this step. We will treat generation as a batch-inference job that we would want to run
ourselves.

\subsection{Teacher generation as a batch job}

\begin{napkinmath}[How big is the generation job?]
12{,}000 problems $\times$ $\approx$3{,}000 output tokens is $\approx 36$M generated tokens (the
prompts add little). For a self-hosted 671B mixture-of-experts teacher with $\approx$37B active
parameters per token, decode costs about $2\times 37\text{B} \approx 74$\,GFLOP per token, i.e.
$\approx 2.7\times10^{18}$\,FLOPs in total---trivial compute for a multi-node serving deployment.
The binding constraint is memory bandwidth and KV capacity, so the job should run at the
largest batch size the deployment supports, with throughput-oriented scheduling
(\cref{ch:serving,ch:efficient}). It has no latency requirement at all.
\end{napkinmath}

A generation script for a large job must be \emph{resumable} (it will be interrupted),
\emph{idempotent} (re-running must not duplicate records), and \emph{robust} to transient
errors (rate limits, timeouts). \Cref{lst:distill:gen} shows the shape of such a script; the
teacher is any callable, e.g.\ a thin client around an OpenAI-compatible endpoint served by
vLLM or SGLang.

\begin{lstlisting}[caption={Resumable, verifiable teacher-trace generation
(\texttt{llmsys/distill.py}).},label={lst:distill:gen}]
def generate_teacher_dataset(problems, teacher_fn, out_path,
                             verify_fn=None, max_retries=2):
    done = set()
    if os.path.exists(out_path):                      # resume: skip finished ids
        with open(out_path) as f:
            done = {json.loads(line)["id"] for line in f if line.strip()}
    stats = dict(skipped_existing=0, written=0, rejected=0, failed=0)
    with open(out_path, "a") as f:                    # append-only JSONL
        for i, prob in enumerate(problems):
            pid = prob.get("id", i)
            if pid in done:
                stats["skipped_existing"] += 1
                continue
            out = None
            for _ in range(max_retries + 1):
                try:
                    out = teacher_fn(render_math_prompt(prob["problem"]))
                    break
                except Exception:                     # rate limit, timeout, ...
                    out = None
            if out is None:
                stats["failed"] += 1
                continue
            if verify_fn and not verify_fn(out["answer"], prob["answer"]):
                stats["rejected"] += 1                # rejection sampling
                continue
            rec = dict(id=pid, problem=prob["problem"], gtruth=prob["answer"], **out)
            f.write(json.dumps(rec) + "\n")
            stats["written"] += 1
    return stats
\end{lstlisting}

\begin{systemslens}[Rejection sampling is cheap quality control]
Passing a \ident{verify_fn} (the verifier of \cref{ch:eval}) keeps only traces whose final
answer is correct. The cost is the fraction of teacher tokens thrown away; the benefit is that
the student never imitates a wrong solution. The book's main experiment trains on all traces
(the teacher is right about 90\% of the time on these problems), so rejection sampling is a
natural ablation. With a per-problem budget one can also sample several traces and keep, say,
the shortest correct one---a first step toward \emph{efficient} reasoning
(\cref{sec:distill:efficient}).
\end{systemslens}

% =====================================================================================
\section{Loading the distillation data}\label{sec:distill:load}\covers{RM 8.3}

Each record of the published dataset contains four fields: the problem statement, the
ground-truth answer, the teacher's reasoning trace, and the teacher's final user-facing
response (which ends in a \texttt{\textbackslash boxed\{\}} answer). Different generators name these fields
differently, so our loader maps them onto one schema---\code{problem}, \code{gtruth},
\code{reasoning}, \code{answer}---via \ident{normalize_record}. Because the file is large
(roughly 100\,MB of JSON for 12k records), the download helper should cache a local copy and
reuse it on later runs; any HTTP client with a fallback mirror will do.

The training target combines the trace and the response. Following the convention of Qwen3
and DeepSeek-R1 reasoning models, we wrap the trace in think tags:

\begin{lstlisting}[caption={Turning a teacher record into the student's target text.},
label={lst:distill:format}]
def format_distilled_answer(record, use_think=True):
    answer = str(record.get("answer", "")).strip()
    if not answer:
        raise ValueError("record has an empty 'answer' field")
    reasoning = str(record.get("reasoning", "")).strip()
    if use_think:
        return f"<think>{reasoning}</think>\n\n{answer}"
    return f"{reasoning}\n\n{answer}" if reasoning else answer
\end{lstlisting}

The tags are not required for distillation to work, but they cleanly separate the part of the
output a user interface may want to hide (the trace) from the part it shows. Since the records
also carry the ground truth, we can measure the \emph{teacher's} accuracy on the training set
with the verifier of \cref{ch:eval}; the source book reports about 90.6\% for DeepSeek-R1 on
these 12k problems (and about 91\% on MATH-500), far above the 0.6B base model's roughly 15\%
on MATH-500. The gap between teacher and student accuracy is the headroom distillation tries
to close.
% =====================================================================================
\section{Building training examples}\label{sec:distill:examples}\covers{RM 8.4}

In RLVR every prompt was used once and its rollouts were discarded, so preprocessing happened
on the fly. Distillation revisits the same examples for several epochs, so it pays to do the
work once: render prompts, format targets, tokenise, and store token ids together with the
prompt length. We proceed in three stages: load the tokenizer, tokenise the whole dataset,
then filter and split it.

\subsection{The tokenizer and the chat format}\label{sec:distill:tok}\covers{RM 8.4.1}

We use the tokenizer of the Qwen3 \emph{reasoning} checkpoint because its vocabulary already
contains \texttt{<think>} and \texttt{</think>} as single special tokens (ids 151667 and
151668); with the base tokenizer the tags would be split into several ordinary tokens, which
works but is wasteful and less clean. The prompt is wrapped in Qwen's ChatML template
(\cref{ch:inference}):
\begin{lstlisting}[style=plain]
<|im_start|>user
Solve the following math problem. ... Problem: ...<|im_end|>
<|im_start|>assistant
\end{lstlisting}
The final line is the \emph{generation prompt}: the model's job begins right after it.
The target answer must \emph{not} be chat-wrapped---otherwise a second
\texttt{<|im\_start|>assistant} header would appear in the middle of the sequence. The full
training sequence is therefore
\[
\underbrace{\text{chat-wrapped prompt}}_{\text{\texttt{prompt\_len} tokens}}\ \Vert\
\underbrace{\texttt{<think>}\ \text{trace}\ \texttt{</think>}\ \text{answer}}_{\text{raw, not wrapped}}\ \Vert\
\texttt{<|im\_end|>},
\]
where the final end-of-turn token teaches the student to stop. Neither the chat template nor
the reasoning tokenizer is strictly necessary, but matching the conventions of Qwen3's own
reasoning models keeps training and evaluation consistent: whichever evaluation script you use
later must load the same tokenizer variant.

\begin{figure}[t]
\centering
\resizebox{\linewidth}{!}{%
\begin{tikzpicture}[tk/.style={draw=white,minimum height=0.6cm,font=\scriptsize,inner sep=3pt,
  anchor=west,outer sep=0pt}]
\node[tk,fill=fdata!35] (p1) at (0,0) {\texttt{<|im\_start|>user}};
\node[tk,fill=fdata!20] (p2) at (p1.east) {instruction + problem};
\node[tk,fill=fdata!35] (p3) at (p2.east) {\texttt{<|im\_end|> <|im\_start|>assistant}};
\node[tk,fill=fcompute!45] (a1) at (p3.east) {\texttt{<think>}};
\node[tk,fill=fcompute!25] (a2) at (a1.east) {teacher reasoning trace};
\node[tk,fill=fcompute!45] (a3) at (a2.east) {\texttt{</think>}};
\node[tk,fill=fcompute!25] (a4) at (a3.east) {answer \texttt{\textbackslash boxed\{6\}}};
\node[tk,fill=fmemory!50] (e) at (a4.east) {\texttt{<|im\_end|>}};
\draw[decorate,decoration={brace,amplitude=5pt,mirror}] (p1.south west) -- (p3.south east)
  node[midway,below=6pt,font=\scriptsize,align=center]{prompt: \texttt{encode(..., chat)}\\ length = \texttt{prompt\_len}; \textbf{no loss}};
\draw[decorate,decoration={brace,amplitude=5pt,mirror}] (a1.south west) -- (a4.south east)
  node[midway,below=6pt,font=\scriptsize,align=center]{target: \texttt{encode(format\_distilled\_answer(r))}\\ not chat-wrapped; \textbf{loss on every token}};
\draw[decorate,decoration={brace,amplitude=5pt,mirror}] (e.south west) -- (e.south east)
  node[midway,below=6pt,font=\scriptsize,align=center]{EOS\\(loss)};
\node[font=\scriptsize,anchor=south west] at ([yshift=2pt]p1.north west) {\texttt{token\_ids = prompt\_ids + answer\_ids + [eos\_id]}};
\end{tikzpicture}}
\caption{Layout of one distillation example. Only the prompt is chat-wrapped. We store the
prompt length so that the loss can be restricted to the teacher's tokens and the final
end-of-turn token.}
\label{fig:distill:layout}
\end{figure}

\subsection{Formatting and tokenising the dataset}\label{sec:distill:build}\covers{RM 8.4.2}

\Cref{lst:distill:build} applies three steps to every record: render and tokenise the prompt,
format and tokenise the target, concatenate both with the end token and remember
\code{prompt\_len}. Records with missing fields are counted and skipped rather than crashing
the job. The encoders are passed in as functions, so the same code works with the Qwen3
tokenizer (in \texttt{run.py}) and with the lab's character tokenizer.

\begin{lstlisting}[caption={Tokenising the dataset once (\texttt{llmsys/distill.py}).},
label={lst:distill:build}]
def build_examples(records, encode_prompt, encode_answer, eos_id,
                   use_think=True, render=render_math_prompt):
    examples, skipped = [], 0
    for rec in records:
        try:
            p = encode_prompt(render(rec["problem"]))           # chat-wrapped
            a = encode_answer(format_distilled_answer(rec, use_think))  # raw
        except (KeyError, TypeError, ValueError):
            skipped += 1
            continue
        ids = p + a + [eos_id]
        if len(ids) < 2 or len(p) == 0:
            skipped += 1
            continue
        examples.append({"token_ids": ids, "prompt_len": len(p)})
    return examples, skipped
\end{lstlisting}

Always decode one example back to text and read it. It should show the chat-wrapped prompt,
the trace inside think tags, the boxed answer, and the end-of-turn token, in that order and
with no duplicated headers.

\subsection{Filtering by length and splitting}\label{sec:distill:filter}\covers{RM 8.4.3}

Reasoning traces are long and heavy-tailed. On the DeepSeek-R1 MATH traces the source book
reports an average of about 2{,}950 tokens per formatted example, a minimum of 236, and a
maximum above 42{,}000. Training cost per example grows at least linearly with length (MLPs,
projections) and partly quadratically (attention), and peak activation memory is set by the
\emph{longest} sequence in a step, so outliers are expensive. The book caps examples at 2{,}048
tokens, which removes about 44\% of the data (6{,}695 of 12{,}000 remain, averaging about 1{,}180
tokens), shuffles with a fixed seed and holds out a deliberately tiny validation set of 25
examples to keep evaluation cheap. Our helpers compute the same quantities
(\cref{lst:distill:filter}).

\begin{lstlisting}[caption={Length statistics, filtering and splitting.},label={lst:distill:filter}]
def length_stats(examples, answer_only=False):
    lens = [len(e["token_ids"]) - (e["prompt_len"] if answer_only else 0)
            for e in examples]
    srt = sorted(lens)
    q = lambda f: srt[min(len(srt) - 1, int(f * len(srt)))]
    return dict(n=len(lens), mean=sum(lens) / len(lens), min=srt[0], max=srt[-1],
                p50=q(0.5), p90=q(0.9), p99=q(0.99),
                argmin=lens.index(srt[0]), argmax=lens.index(srt[-1]))

def filter_by_max_len(examples, max_len):
    return [e for e in examples if len(e["token_ids"]) <= max_len]

def split_train_val(examples, val_size, seed=123):
    ex = list(examples)
    random.Random(seed).shuffle(ex)
    return ex[val_size:], ex[:val_size]
\end{lstlisting}

\begin{figure}[t]
\centering
\begin{tikzpicture}
\begin{axis}[width=0.85\linewidth,height=5cm,xlabel={tokens per example},
  ylabel={density (illustrative)},xmin=0,xmax=10000,ymin=0,
  domain=50:10000,samples=200,yticklabels={},grid=major,grid style={gray!15},
  scaled x ticks=false,xticklabel style={/pgf/number format/1000 sep={,}}]
% log-normal with median ~2200, sigma 0.75
\addplot[thick,sysblue,fill=sysblue!15] {exp(-(ln(x)-ln(2200))^2/(2*0.75^2))/(x*0.75)} \closedcycle;
\addplot[very thick,warnred] coordinates {(2048,0) (2048,0.00025)};
\node[font=\scriptsize,warnred,anchor=west,align=left] at (axis cs:2150,0.00022) {cap at 2{,}048:\\keep $\approx$56\%};
\node[font=\scriptsize,anchor=west,align=left] at (axis cs:5500,0.00008) {long tail to $>$40k tokens\\(clipped from plot)};
\end{axis}
\end{tikzpicture}
\caption{Schematic length distribution of reasoning traces (a log-normal shape, not the
measured histogram) and the effect of a length cap. A cap is the single most effective knob
for fitting distillation into a small GPU, but it biases the data (see the pitfall below).}
\label{fig:distill:lengths}
\end{figure}

\begin{pitfall}[A length cap is a difficulty filter]
Long traces are disproportionately the solutions to hard problems (and the teacher's failed,
meandering attempts). Dropping everything above 2{,}048 tokens therefore removes most of the
hardest problems and teaches the student a style that rarely needs long reasoning. That is
sometimes desirable (short answers are cheap to serve) and sometimes the reason a distilled
model plateaus. Inspect accuracy by MATH difficulty level before and after filtering, and
prefer \emph{packing} and \emph{activation checkpointing} (\cref{sec:distill:longseq}) to
raise the cap when memory allows.
\end{pitfall}

A natural check, and the subject of an exercise, is whether the random split produced train and
validation sets with similar length distributions; with only 25 validation examples it may not,
and validation loss is length-sensitive.

% =====================================================================================
\section{Loading the student}\label{sec:distill:model}\covers{RM 8.5}

We start from the pretrained Qwen3-0.6B \emph{base} model rather than from an RL-trained
checkpoint. That keeps the experiment clean---any gain is attributable to the distilled
traces---and mirrors the common practical setting of adapting a general base model with
teacher data. The base model's own tokenizer is not used; the reasoning tokenizer from
\cref{sec:distill:tok} has the same vocabulary plus the think tags already defined as special
tokens, and the embedding matrix already has rows for them. With our core code
(\cref{ch:modern}), loading is one line:
\begin{lstlisting}
model = load_qwen3_from_dir("qwen3-base", "0.6B", dtype=torch.bfloat16, device=device)
\end{lstlisting}

\begin{napkinmath}[Will it fit?]
Qwen3-0.6B has $\approx$0.60B parameters (with tied embeddings). Full fine-tuning with AdamW in
mixed precision costs $\approx 16$\,B/param, i.e.\ $\approx$9.5\,GB for weights, gradients and
optimizer state. Activations for one 2{,}048-token sequence add a few GB; the logits alone are
$2048\times151{,}936\times 4$\,B $\approx 1.2$\,GB in fp32, plus the same again for their gradient.
The source book's measured peak of about 15\,GB at batch size 1 is consistent with this budget.
Doubling the cap to 4{,}096 tokens roughly doubles activation and logit memory---this is where
the techniques of \cref{sec:distill:longseq} pay off.
\end{napkinmath}

% =====================================================================================
\section{The training and validation losses}\label{sec:distill:loss}\covers{RM 8.6}

Distillation's loss is the familiar next-token cross-entropy of \cref{ch:pretrain}, restricted to
the target part of each sequence. It is worth connecting it explicitly to the log-probabilities
of \cref{ch:refine,ch:grpo}. For an example with prompt $x$ and target $y=(y_1,\dots,y_n)$,
\begin{equation}
\mathcal{L}_{\text{CE}} = -\frac{1}{n}\sum_{t=1}^{n}\log p_\theta(y_t\mid x, y_{<t})
= -\frac{1}{n}\log p_\theta(y\mid x),
\label{eq:distill:ce}
\end{equation}
the \emph{negative average log-probability} of the teacher's tokens. Summed log-probabilities
grow with length and cannot be compared across examples; dividing by $n$ gives a per-token
quantity. If a model assigns joint log-probability $-16.6$ to a five-token continuation, its
cross-entropy on that continuation is $16.6/5 = 3.32$ nats per token.

\begin{figure}[t]
\centering
\resizebox{\linewidth}{!}{%
\begin{tikzpicture}[c/.style={draw,minimum width=1.05cm,minimum height=0.55cm,font=\scriptsize,inner sep=1pt},x=1.15cm]
\node[font=\scriptsize\bfseries,anchor=east] at (-0.3,1.5) {sequence};
\node[font=\scriptsize\bfseries,anchor=east] at (-0.3,0.6) {model input};
\node[font=\scriptsize\bfseries,anchor=east] at (-0.3,-0.3) {logit predicts};
\node[font=\scriptsize\bfseries,anchor=east] at (-0.3,-1.2) {in loss?};
\foreach \i/\t/\c in {0/$p_0$/fdata,1/$p_1$/fdata,2/$p_2$/fdata,3/$a_0$/fcompute,4/$a_1$/fcompute,5/$a_2$/fcompute,6/\textsc{eos}/fmemory}
  \node[c,fill=\c!30] at (\i,1.5) {\t};
\foreach \i/\t/\c in {0/$p_0$/fdata,1/$p_1$/fdata,2/$p_2$/fdata,3/$a_0$/fcompute,4/$a_1$/fcompute,5/$a_2$/fcompute}
  \node[c,fill=\c!30] at (\i,0.6) {\t};
\foreach \i/\t in {0/$p_1$,1/$p_2$,2/$a_0$,3/$a_1$,4/$a_2$,5/\textsc{eos}}
  \node[c] at (\i,-0.3) {\t};
\foreach \i/\t/\c in {0/no/gray,1/no/gray,2/yes/keygreen,3/yes/keygreen,4/yes/keygreen,5/yes/keygreen}
  \node[c,fill=\c!25] at (\i,-1.2) {\t};
\draw[warnred,very thick,-{Stealth}] (2,-1.85) -- (2,-1.55);
\node[font=\scriptsize,warnred,anchor=north] at (2,-1.85) {index \texttt{prompt\_len - 1}};
\node[font=\scriptsize,align=left,anchor=west] at (7.0,0.1) {$T$ tokens in, $T-1$ inputs;\\logit at position $t$ is scored\\against token $t+1$.\\Prompt length here $=3$.};
\end{tikzpicture}}
\caption{Answer-only loss alignment. The model sees all tokens but the last; the logit at
position $t$ predicts token $t+1$. The first answer token $a_0$ (at index \texttt{prompt\_len})
is therefore predicted by the logit at index \texttt{prompt\_len}$-1$, which is where the loss
starts.}
\label{fig:distill:shift}
\end{figure}

We compute the loss only over the answer: the prompt is context the student is given, not
something it should be penalised for failing to predict (\cref{fig:distill:shift}). Because of the
one-token shift, the slice starts at logit index \code{prompt\_len - 1}:

\begin{lstlisting}[caption={Answer-only cross-entropy for one example.},label={lst:distill:loss}]
def answer_only_loss(model, example, device="cpu"):
    ids = torch.tensor(example["token_ids"], dtype=torch.long, device=device)
    logits = model(ids[:-1][None])[0]           # (T-1, V)
    start = max(example["prompt_len"] - 1, 0)   # logit that predicts answer[0]
    return F.cross_entropy(logits[start:].float(), ids[1:][start:])
\end{lstlisting}

The lab verifies \cref{eq:distill:ce} numerically: \ident{answer_only_loss} equals
$-\ident{sequence_logprob}/n$ computed with the generic helper of \cref{ch:refine}. In
practice \texttt{F.cross\_entropy} is preferred because it fuses log-softmax and gather and
never materialises the full log-probability tensor separately.

\begin{pitfall}[Off by one]
Starting the slice at \code{prompt\_len} instead of \code{prompt\_len - 1} silently drops the
first answer token from the loss---which, with think tags, is the \texttt{<think>} token the
student must learn to emit first. The model trains fine and then forgets to open its
reasoning at inference time. Test the alignment on a toy example where you can enumerate every
position.
\end{pitfall}

\paragraph{Validation loss.} A helper averages \ident{answer_only_loss} over a list of examples
in \texttt{eval} mode under \texttt{torch.no\_grad()}, restoring training mode afterwards so it
is safe to call inside the loop. Training loss is measured on the example just optimised and is
noisy; validation loss on a fixed held-out set is the cleaner signal of generalisation.

\subsection{Batching the loss correctly}

The book's loop processes one example per step to keep memory low. Batching several examples
requires padding (\cref{ch:instruct}) and a decision about normalisation that is easy to get
wrong. With right padding and $-100$ (ignored) targets on prompt and pad positions, there are
two reasonable reductions:
\begin{align}
\mathcal{L}_{\text{token}} &= \frac{\sum_i \sum_{t\in\mathcal{A}_i} \ell_{i,t}}{\sum_i |\mathcal{A}_i|},&
\mathcal{L}_{\text{seq}} &= \frac{1}{B}\sum_i \frac{1}{|\mathcal{A}_i|}\sum_{t\in\mathcal{A}_i}\ell_{i,t}.
\end{align}
$\mathcal{L}_{\text{token}}$ (what cross-entropy with \code{ignore\_index=-100} gives on the
flattened batch) weights each \emph{token} equally, so long traces dominate the gradient;
$\mathcal{L}_{\text{seq}}$ weights each \emph{example} equally and reproduces the average of
per-example losses. Neither is wrong, but switching between them changes the effective learning
rate per example by up to the ratio of the longest to shortest trace---and they interact with
gradient accumulation exactly as described in \cref{sec:rlsys:packing}.

\begin{lstlisting}[caption={Padded batch with prompt and pad positions masked out.},
label={lst:distill:collate}]
def collate_answer_only(batch, pad_id, device="cpu"):
    T = max(len(e["token_ids"]) for e in batch)
    inp = torch.full((len(batch), T - 1), pad_id, dtype=torch.long)
    tgt = torch.full((len(batch), T - 1), IGNORE, dtype=torch.long)
    for i, e in enumerate(batch):
        ids = torch.tensor(e["token_ids"], dtype=torch.long)
        n = len(ids)
        inp[i, : n - 1] = ids[:-1]
        tgt[i, : n - 1] = ids[1:]
        tgt[i, : max(e["prompt_len"] - 1, 0)] = IGNORE   # no loss on the prompt
    return inp.to(device), tgt.to(device), int((tgt != IGNORE).sum())
\end{lstlisting}

Right padding is safe for training a causal model: real tokens never attend to later
positions, so pads cannot leak into real positions' outputs.

% =====================================================================================
\section{The distillation training loop}\label{sec:distill:loop}\covers{RM 8.7}

The loop (\cref{fig:distill:loop}) is the GRPO loop of \cref{ch:grpo} with the policy-gradient
objective replaced by cross-entropy and with epochs over a fixed dataset instead of fresh
rollouts every step.

\begin{figure}[t]
\centering
\resizebox{0.95\linewidth}{!}{%
\begin{tikzpicture}[st/.style={draw=sysblue,fill=sysbluebg,rounded corners=2pt,font=\scriptsize,
  minimum height=0.75cm,minimum width=2.25cm,align=center},
  opt/.style={st,draw=srcgray,fill=gray!8,dashed},
  arr/.style={-{Stealth[length=2mm]},thick,sysblue!70}]
\node[st] (a) at (0,0) {1\ \ model + AdamW};
\node[st] (b) at (2.8,0) {2\ \ for each epoch};
\node[st] (c) at (5.6,0) {3\ \ shuffle (seeded)};
\node[st] (d) at (8.4,0) {4\ \ for each example\\or batch};
\node[st] (e) at (8.4,-1.4) {5\ \ zero grads};
\node[st] (f) at (5.6,-1.4) {6\ \ answer-only CE};
\node[st] (g) at (2.8,-1.4) {7\ \ backward};
\node[st] (h) at (0,-1.4) {8\ \ clip, step};
\node[opt] (i) at (0,-2.8) {9\ \ every $k$ steps:\\val loss, CSV log};
\node[opt] (j) at (2.8,-2.8) {10\ \ end of epoch:\\checkpoint};
\draw[arr] (a)--(b); \draw[arr] (b)--(c); \draw[arr] (c)--(d); \draw[arr] (d)--(e);
\draw[arr] (e)--(f); \draw[arr] (f)--(g); \draw[arr] (g)--(h); \draw[arr] (h)--(i);
\draw[arr] (i)--(j);
\draw[arr,dashed] (h.west) -- ++(-0.45,0) |- node[pos=0.25,left,font=\scriptsize]{next} ([yshift=0.5cm]a.north) -| (d.north);
\end{tikzpicture}}
\caption{The distillation training loop. Steps 5--8 are standard supervised fine-tuning;
the dashed arrow returns to step 4 for the next example; steps 9--10 are optional bookkeeping
(validation loss, metrics file, checkpoints).}
\label{fig:distill:loop}
\end{figure}

\begin{lstlisting}[caption={Training loop for hard distillation (\texttt{llmsys/distill.py},
abridged).},label={lst:distill:train}]
def train_distillation(model, train_examples, val_examples, device="cpu", epochs=2,
                       lr=5e-6, grad_clip=1.0, seed=123, log_every=50,
                       batch_size=1, pad_id=0, checkpoint_dir=None, csv_path=None):
    opt = torch.optim.AdamW(model.parameters(), lr=lr)
    rng, step, rows = random.Random(seed), 0, []
    model.train()
    for epoch in range(1, epochs + 1):
        order = list(train_examples)
        rng.shuffle(order)                                  # new order each epoch
        for s in range(0, len(order), batch_size):
            batch = order[s:s + batch_size]
            step += 1
            opt.zero_grad(set_to_none=True)
            loss = (answer_only_loss(model, batch[0], device) if batch_size == 1
                    else batch_answer_only_loss(model, batch, pad_id, device))
            loss.backward()
            if grad_clip is not None:
                torch.nn.utils.clip_grad_norm_(model.parameters(), grad_clip)
            opt.step()
            if log_every and step % log_every == 0:
                val = evaluate_examples(model, val_examples, device)
                row = dict(epoch=epoch, step=step, train_loss=loss.item(), val_loss=val)
                rows.append(row)
                if csv_path:
                    append_csv(csv_path, row)
        if checkpoint_dir:                                  # one checkpoint per epoch
            torch.save(model.state_dict(),
                       f"{checkpoint_dir}/distill-step{step:06d}-epoch{epoch}.pth")
    return rows
\end{lstlisting}

The hyperparameters are typical for full fine-tuning of a small pretrained model: a learning
rate around $5\times10^{-6}$ to $10^{-5}$ with AdamW, gradient clipping at norm 1.0 (as in the
GRPO chapters) to tame occasional large gradients from unusual examples, and a validation
interval chosen to balance feedback against evaluation overhead. A smoke test on ten examples
for two epochs should show both losses falling; a real run uses the whole filtered set for
several epochs. Writing metrics to a CSV file after every evaluation makes the run inspectable
while it is still going and plottable afterwards.

\begin{systemslens}[Batch size 1 is a memory decision, not a throughput one]
At batch size 1 every step is a single sequence of $\sim$1{,}200 tokens; for a 0.6B model on a
modern GPU that leaves the matrix units badly underused, and the optimizer step (which touches
every parameter) is amortised over very few tokens. Throughput improves with (i) gradient
accumulation to reach a larger effective batch, (ii) padded mini-batches bucketed by length,
or best (iii) \emph{packed} rows with block-diagonal attention (\cref{sec:rlsys:packing}), which
in the lab reach $>$95\% real tokens versus $\sim$40--70\% for padding on heavy-tailed lengths.
Remember to keep the loss normalisation fixed when you change batching.
\end{systemslens}

% =====================================================================================
\section{Evaluating the distilled model}\label{sec:distill:eval}\covers{RM 8.8}

Evaluation has two layers. The \emph{learning curves}---training and validation loss over steps,
plotted from the CSV log with a secondary axis marking epochs---tell us whether optimisation is
healthy. In the book's three-epoch run on the DeepSeek-R1 traces the validation loss drops
steeply at first and then flattens, the usual sign that further epochs give diminishing
returns. The training curve is noisier because each point is a single example.

The second layer is \emph{task accuracy} on MATH-500 with the verifier of \cref{ch:eval},
evaluating each epoch's checkpoint with the reasoning tokenizer and a generous generation budget
(the book uses 4{,}096 new tokens). The source book reports the following (we quote its numbers;
they were not reproduced here):
\begin{center}\small
\begin{tabular}{@{}lccc@{}}
\toprule
Model & Epoch & Final val.\ loss & MATH-500\\
\midrule
Qwen3-0.6B base (no distillation) & -- & -- & 15.2\%\\
Qwen3-0.6B official reasoning model & -- & -- & 48.2\%\\
Distilled from DeepSeek-R1 traces & 1 / 2 / 3 & 0.544 / 0.535 / 0.534 & 30.6 / 32.4 / 33.6\%\\
Distilled from Qwen3-235B-A22B traces & 1 / 2 / 3 & 0.404 / 0.396 / 0.395 & 45.0 / 43.8 / 44.2\%\\
\bottomrule
\end{tabular}
\end{center}

Three observations. First, a few hours of supervised training on $\sim$6.7k filtered traces
roughly doubles to triples the base model's accuracy. Second, the students remain far below
their teachers (both above 90\% on MATH-500): a 0.6B model cannot absorb everything a 671B or
235B model knows, and larger students (4B, 8B) would close more of the gap at higher training
cost. Third, the \emph{same-family} teacher works markedly better: Qwen3-235B-A22B traces yield a
lower validation loss and a student nearly matching Qwen's own distilled 0.6B model, which was
trained on far more data. A plausible explanation is closer alignment of tokenisation, formatting
and style, which makes the targets easier for the student to imitate---a reminder that ``teacher
quality'' includes how learnable its outputs are, not only how correct.

\begin{keyidea}[Distillation in one sentence]
Distillation is supervised fine-tuning on teacher-written reasoning, with an answer-only
cross-entropy that equals the negative mean log-probability of the teacher's tokens; it buys
a large fraction of the teacher's reasoning ability at a small fraction of RL's cost, bounded by
the student's capacity and by how imitable the teacher is.
\end{keyidea}

A natural ablation (Exercise~\ref{ex:distill:nothink} and the lab) asks how much the explicit
think tags matter. In our synthetic lab task we can go further and ask how much the trace
itself matters, by training students on the full trace, on the trace without tags, and on the
answer alone.

\begin{figure}[t]
\centering
\begin{tikzpicture}
\begin{axis}[ybar,width=0.88\linewidth,height=5.4cm,bar width=11pt,ymin=0,ymax=100,
  symbolic x coords={think 6ep,plain 6ep,direct 6ep,direct 18ep},xtick=data,
  xticklabels={{trace + tags\\6 epochs},{trace, no tags\\6 epochs},{answer only\\6 epochs},{answer only\\18 epochs}},
  xticklabel style={align=center,font=\scriptsize},
  ylabel={held-out accuracy (\%)},enlarge x limits=0.15,
  legend style={font=\scriptsize,at={(0.5,1.03)},anchor=south,legend columns=2},
  nodes near coords,nodes near coords style={font=\tiny},grid=major,grid style={gray!15}]
\addplot[fill=fcompute!60] coordinates {(think 6ep,86) (plain 6ep,85) (direct 6ep,1) (direct 18ep,85)};
\addplot[fill=fmemory!70] coordinates {(think 6ep,61) (direct 18ep,1)};
\legend{2-digit addition,3-digit addition}
\end{axis}
\end{tikzpicture}
\caption{Lab ablation on the synthetic addition task (measured on CPU in this environment with
\texttt{bench.py}, seed 0, 1{,}500 training problems, 100 unseen test problems, a 4-layer
$d=128$ student): held-out exact-match accuracy of students trained on the teacher's full trace
with think tags, on the trace without tags, and on the final answer only. Missing bars were not
run.}
\label{fig:distill:labresult}
\end{figure}

The result (\cref{fig:distill:labresult}) is instructive. With the same 1{,}500 problems and six
epochs, students trained on the trace reach 86\% (with tags) and 85\% (without) exact-match
accuracy on unseen 2-digit problems, while the answer-only student reaches 1\%. Training the
answer-only student three times longer---roughly equalising the number of target tokens seen---lets
it catch up on 2-digit problems (85\%), but on 3-digit problems the trace-trained student reaches
61\% after six epochs while the answer-only student is still at 1\% after eighteen. A second seed
of the 2-digit run gave 72\% (trace) and 4\% (answer only), so the trace effect is robust while the
exact numbers vary by seed. The trace makes the task \emph{learnable} for a small model by turning
one hard mapping into a sequence of easy local steps---the same intuition behind distilling
reasoning traces rather than answers. The price appears at inference: the traced student emits
about 2.7--3.3$\times$ more tokens per answer (31.5 vs.\ 11.5 target tokens for 2 digits, 41.0 vs.\
12.6 for 3 digits), which is exactly the cost that \cref{sec:distill:efficient} tries to reduce.
% =====================================================================================
\section{Soft distillation as a systems problem}\label{sec:distill:soft}

When teacher and student share a tokenizer---as with the Qwen3 family---soft distillation
becomes possible and often more sample-efficient. Its difficulties are almost entirely systems
difficulties: where the teacher's distributions come from, how they are stored, and how the
loss is computed without materialising enormous tensors.

\subsection{Forward and reverse KL}

\Cref{eq:distill:kd} uses the \emph{forward} KL, $\KL(p_T\|p_S)$, which penalises the student
for putting low probability where the teacher puts mass; it is \emph{mass-covering}, pushing the
student to spread probability over every mode the teacher considers plausible. The
\emph{reverse} KL, $\KL(p_S\|p_T)$, penalises the student for putting mass where the teacher
does not; it is \emph{mode-seeking}. A small student that cannot represent all of a teacher's
modes behaves very differently under the two (\cref{fig:distill:kl}): forward KL produces a broad
compromise that may assign probability to regions neither mode supports, while reverse KL
commits to one mode. MiniLLM (Gu et al., 2024) argued for reverse KL when distilling generative
LLMs, and on-policy methods such as GKD (Agarwal et al., 2024) compute the divergence on
sequences sampled \emph{from the student}, which fixes the train/test mismatch of training only
on teacher-written prefixes.

\begin{figure}[t]
\centering
\begin{tikzpicture}
\begin{axis}[width=0.85\linewidth,height=4.6cm,xmin=-6,xmax=6,ymin=0,ymax=0.45,
  domain=-6:6,samples=200,yticklabels={},xticklabels={},axis lines=left,
  legend style={font=\scriptsize,at={(0.99,0.98)},anchor=north east}]
\addplot[thick,black,fill=gray!15] {0.5*exp(-(x+2.2)^2/(2*0.6^2))/(0.6*sqrt(2*pi)) + 0.5*exp(-(x-2.2)^2/(2*0.6^2))/(0.6*sqrt(2*pi))} \closedcycle;
\addlegendentry{teacher $p_T$ (bimodal)}
\addplot[very thick,sysblue,dashed] {exp(-(x)^2/(2*2.3^2))/(2.3*sqrt(2*pi))};
\addlegendentry{student, forward KL (covers both)}
\addplot[very thick,warnred] {exp(-(x-2.2)^2/(2*0.65^2))/(0.65*sqrt(2*pi))};
\addlegendentry{student, reverse KL (picks one)}
\end{axis}
\end{tikzpicture}
\caption{A unimodal student fitted to a bimodal teacher. Forward KL (classic KD) spreads mass
over both modes and the valley between them; reverse KL (MiniLLM, on-policy distillation) locks
onto one mode. Curves are schematic.}
\label{fig:distill:kl}
\end{figure}

\begin{lstlisting}[caption={Token-level KD loss with temperature and direction, and a sparse
top-$k$ variant (\texttt{llmsys/distill.py}).},label={lst:distill:kd}]
def kd_loss(student_logits, teacher_logits, mask=None, temperature=1.0,
            direction="forward"):
    T = temperature
    ls = F.log_softmax(student_logits.float() / T, dim=-1)
    lt = F.log_softmax(teacher_logits.float() / T, dim=-1)
    if direction == "forward":
        kl = (lt.exp() * (lt - ls)).sum(-1)        # KL(p_T || p_S)
    else:
        kl = (ls.exp() * (ls - lt)).sum(-1)        # KL(p_S || p_T)
    m = torch.ones_like(kl) if mask is None else mask.float()
    return (kl * m).sum() / m.sum().clamp(min=1) * T * T

def topk_teacher(teacher_logits, k):
    lp = F.log_softmax(teacher_logits.float(), dim=-1)
    vals, idx = torch.topk(lp, k, dim=-1)
    return vals.to(torch.float16), idx.to(torch.int32)   # 6 bytes per entry

def sparse_kd_loss(student_logits, topk_logp, topk_idx, mask=None):
    ls_k = F.log_softmax(student_logits.float(), -1).gather(-1, topk_idx.long())
    pt = torch.softmax(topk_logp.float(), dim=-1)          # renormalise over k
    kl = (pt * (pt.clamp_min(1e-12).log() - ls_k)).sum(-1)
    m = torch.ones_like(kl) if mask is None else mask.float()
    return (kl * m).sum() / m.sum().clamp(min=1)
\end{lstlisting}

\subsection{Offline, online, and sparse teachers}

There are three ways to supply $p_T$:
\begin{description}
\item[Offline dense.] Precompute and store full distributions. As the napkin math in
\cref{sec:distill:intro} showed, this is impractical beyond toy scales.
\item[Offline sparse.] Store the top-$k$ log-probabilities and token ids per position
($k=8$--$64$), renormalise over the support at training time (\ident{sparse_kd_loss}). For
peaked reasoning distributions most mass sits in a handful of tokens, so small $k$ loses little;
the lab checks that $k=V$ recovers the dense loss exactly (up to fp16 storage).
\item[Online.] Run the teacher's forward pass inside the training step. No storage, exact
distributions, and on-policy variants become possible---but every training step now pays for a
teacher forward pass (for a teacher $10\times$ the student's size, more than the student's own
forward and backward together), and both models must fit in memory or be served separately.
\end{description}

\begin{napkinmath}[Online teacher cost]
Training the student costs $\approx 6N_S$ FLOPs per token; an online teacher forward pass costs
$\approx 2N_T$. With $N_S=0.6$B and an 8B teacher, the teacher adds $16/3.6\approx 4.4\times$ the
student's training compute. With a 235B-A22B mixture-of-experts teacher it adds
$2\times 22/3.6\approx 12\times$ (counting active parameters) and needs a multi-GPU serving
deployment next to the trainer---the same placement questions as RL (\cref{sec:rlsys:placement}).
\end{napkinmath}

% =====================================================================================
\section{Training on long traces}\label{sec:distill:longseq}

Raising the 2{,}048-token cap is the most direct way to keep hard problems in the data
(\cref{sec:distill:filter}). Three memory consumers grow with sequence length $T$:

\paragraph{Activations.} Per layer roughly linear in $T$ with FlashAttention-style kernels
(no $T^2$ score matrix stored). Activation checkpointing (\cref{ch:pretrain}) trades a second
forward pass for storing only layer inputs.

\paragraph{Logits.} The output layer produces $T\times V$ values; with $V\approx 152$k the logits
are the \emph{largest single activation} in a small model. At $T=16{,}384$ in fp32 that is
$\approx 10$\,GB, plus a gradient of the same size. The fix is to never materialise them all:
compute the loss in chunks of positions, so that only $\text{chunk}\times V$ logits exist at a
time. \Cref{lst:distill:chunkce} is a minimal version; fused kernels (e.g.\ Liger Kernel's fused
linear cross-entropy) additionally compute the gradient inside the chunk loop so that the chunk
need not be kept for backward.

\begin{lstlisting}[caption={Chunked cross-entropy: peak logit memory is
\texttt{chunk}$\times V$ instead of $T\times V$.},label={lst:distill:chunkce}]
def chunked_cross_entropy(hidden, out_head, targets, chunk=1024):
    total = hidden.new_zeros((), dtype=torch.float32)
    count = (targets != IGNORE).sum().clamp(min=1)
    for s in range(0, hidden.shape[0], chunk):
        logits = out_head(hidden[s:s + chunk]).float()       # (chunk, V)
        total = total + F.cross_entropy(logits, targets[s:s + chunk],
                                        ignore_index=IGNORE, reduction="sum")
    return total / count
\end{lstlisting}

Note that with autograd, \cref{lst:distill:chunkce} still keeps each chunk's logits alive for the
backward pass unless each chunk is wrapped in activation checkpointing; the chunking reduces the
\emph{transient} peak (softmax intermediates, fp32 copies) and is the stepping stone to a fused
kernel.

\paragraph{Attention across devices.} Beyond what one GPU can hold, the sequence itself is
split: context parallelism (ring attention) shards the sequence across GPUs and circulates
key/value blocks (\cref{ch:distributed}). For a 0.6B student this is rarely needed below
$\sim$32k tokens; for larger students on 64k-token traces it is standard.

\begin{figure}[t]
\centering
\begin{tikzpicture}
\begin{axis}[width=0.8\linewidth,height=5cm,xlabel={sequence length $T$ (tokens)},
  ylabel={GB},xmode=log,log basis x=2,xmin=1024,xmax=32768,ymin=0,
  xtick={1024,2048,4096,8192,16384,32768},
  xticklabels={1k,2k,4k,8k,16k,32k},
  legend style={font=\scriptsize,at={(0.02,0.98)},anchor=north west},grid=major,grid style={gray!15}]
\addplot[thick,warnred,mark=*] coordinates {(1024,0.622) (2048,1.245) (4096,2.489) (8192,4.978) (16384,9.957) (32768,19.91)};
\addlegendentry{full fp32 logits $T\cdot V\cdot 4$\,B}
\addplot[thick,sysblue,mark=square*] coordinates {(1024,0.622) (2048,0.622) (4096,0.622) (8192,0.622) (16384,0.622) (32768,0.622)};
\addlegendentry{chunked, chunk $=1024$}
\addplot[thick,dashed,black] coordinates {(1024,9.54) (32768,9.54)};
\addlegendentry{weights+grads+Adam (16 B/param, 0.6B)}
\end{axis}
\end{tikzpicture}
\caption{Memory of the output logits for Qwen3's 151{,}936-token vocabulary versus sequence
length (computed, not measured). Beyond $\sim$16k tokens the logits alone exceed the entire
optimizer state of the 0.6B student; chunking caps them.}
\label{fig:distill:logitmem}
\end{figure}

\paragraph{Packing.} Short traces packed into long rows (block-diagonal attention,
\cref{sec:rlsys:packing}) remove padding; combined with a token-budget batch size this makes step
time and memory nearly independent of the length mix.

% =====================================================================================
\section{Efficient reasoning: making the student cheaper to run}\label{sec:distill:efficient}

A distilled 0.6B model is cheap per token, but the serving cost of a reasoning model is
$\text{(cost per token)} \times \text{(tokens per answer)}$, and traces of thousands of tokens can
make a small model more expensive per query than a larger model that answers directly. Several
levers act on the second factor:
\begin{itemize}
\item \textbf{Choose shorter traces.} When generating the dataset, sample several traces per
problem and keep the shortest \emph{correct} one; or filter by length as in
\cref{sec:distill:filter}, accepting the difficulty bias. The student inherits the teacher's
verbosity, so curating for concision is the cheapest compression.
\item \textbf{Budget control at inference.} ``Budget forcing'' (Muennighoff et al., 2025, the
\emph{s1} paper) caps thinking by forcibly closing the think block, or extends it by appending a
token like ``Wait''; production systems expose a reasoning-effort setting or a thinking-token
budget and let the application choose.
\item \textbf{Length-aware RL after distillation.} A brief RLVR phase (\cref{ch:grpo}) with a
length penalty, or with DAPO's overlong shaping (\cref{ch:grpo2}), can shorten traces while
preserving accuracy; distill-then-RL is a common two-stage recipe.
\item \textbf{Hybrid thinking modes.} Train the student on both thinking and non-thinking
targets (Qwen3's empty \texttt{<think></think>} convention, \cref{ch:modern}) so the caller can
switch reasoning off for easy queries.
\item \textbf{Faster decoding.} A distilled student is also a natural \emph{draft} model for
speculative decoding of its teacher (\cref{ch:efficient}), since distillation increases the
draft--target acceptance rate.
\end{itemize}

\begin{systemslens}[Cost per solved problem]
The metric that matters for deployment is cost per \emph{correct} answer:
$\text{tokens per answer}\times\text{cost per token}/\text{accuracy}$. A student at 33\% accuracy
using 2{,}000 tokens per answer spends about $2000/0.33\approx 6{,}100$ tokens per solved problem;
a length-penalised version that reaches 31\% with 800 tokens spends about $2{,}600$, i.e.\ it is
$2.3\times$ cheaper per solved problem despite being ``worse''. Report both axes.
\end{systemslens}

% =====================================================================================
\section{Future directions for reasoning models}\label{sec:distill:future}\covers{RM 8.9}

The broad recipe seems likely to persist for a while: train the strongest possible reasoning
teacher, collect high-quality traces, and distill them into smaller, cheaper models. Four
directions stand out.
\begin{enumerate}
\item \textbf{Refining the R1-style pipeline.} RLVR (\cref{ch:grpo,ch:grpo2}) for flagship
models, distillation for small ones, and continual improvements to both---better GRPO variants,
better data curation, better teachers.
\item \textbf{Adaptive inference-time compute.} Not every question deserves a long trace. The
surrounding system can decide when to answer directly, when to think longer, and when to stop
early; commercial products in 2025 began exposing automatic reasoning-effort modes of this kind.
This is a systems problem as much as a modelling one (\cref{ch:scaling}).
\item \textbf{Richer rewards.} Most RLVR rewards check only the final answer. Rewards that
judge intermediate steps or the whole solution (process supervision, model-based proof
grading---DeepSeekMath-V2 is a recent example of judging complete solutions during training) may
teach more reliable reasoning, at the cost of building and serving those judges.
\item \textbf{Agents.} Reasoning models are increasingly the core of agent systems that plan,
call tools, and recover from errors over long workflows. This pushes reward design beyond
correctness toward tool-use success, retrieval quality and policy compliance---\emph{multi-reward}
training---and pushes infrastructure toward the agentic rollouts of \cref{sec:rlsys:agentic}.
Distillation remains the practical way to move such behaviours from expensive teachers into
models that can run locally or cheaply.
\end{enumerate}

% =====================================================================================
\section{Conclusions}\label{sec:distill:conclusions}\covers{RM 8.10}

\subsection{What's next}\covers{RM 8.10.1}

The techniques of Part~VI are most interesting in combination: distill a student from a strong
teacher (this chapter), continue with RLVR (\cref{ch:grpo,ch:grpo2}) using the infrastructure of
\cref{ch:rlsys}, and apply self-consistency or self-refinement at deployment time
(\cref{ch:scaling,ch:refine}), all measured with the evaluation harness of \cref{ch:eval} and served
with the engine of \cref{ch:serving}. Running such end-to-end comparisons on your own tasks is
the fastest way to learn which method helps where. The appendices extend the systems side:
Triton kernels (\cref{app:triton}) and lab setup at scale (\cref{app:setup}).

\subsection{Staying current}\covers{RM 8.10.2}

Specific algorithms and benchmarks in this field change monthly; the fundamentals in this book
change slowly: how losses and rewards are computed, the difference between training-time and
inference-time methods, careful evaluation, and the cost models that govern training and
serving. When a new method appears, map it back to these building blocks---it is usually a
recombination. Practical sources: skim new arXiv listings for themes rather than trying to read
everything; read model cards and technical reports, which often contain the most useful details
on prompting conventions and limitations; follow practitioner communities where replications and
failure cases appear early; and use AI research assistants as filters, not as substitutes for
reading the papers. For the systems side, read the design docs and issue trackers of the open
frameworks (\cref{tab:rlsys:frameworks}) and inference engines---that is where today's
bottlenecks are discussed first.

% =====================================================================================
\section{Summary}\label{sec:distill:summary}
\begin{itemize}
\item Distillation trains a small student on outputs of a large teacher. Hard distillation is
SFT on teacher text; soft distillation matches teacher distributions and needs logits and a
shared vocabulary.
\item The teacher runs once, offline, as a resumable batch job; rejection sampling with a
verifier is cheap quality control.
\item Each example is the chat-wrapped prompt plus the (think-tagged) teacher trace and answer
plus an end token; tokenise once, store \code{prompt\_len}, filter by length (a difficulty
bias), split with a fixed seed.
\item The loss is answer-only cross-entropy, equal to the negative mean log-probability of the
teacher's tokens; mind the one-position shift and the token- vs.\ sequence-level normalisation
when batching.
\item A standard SFT loop with AdamW, clipping, validation loss and per-epoch checkpoints
suffices; validation loss and benchmark accuracy are the signals to watch. Same-family teachers
are easier to imitate.
\item Soft distillation is a storage/compute problem: use top-$k$ sparse targets or an online
teacher; forward KL covers modes, reverse KL seeks one.
\item Long traces are limited by activations and, for large vocabularies, by logits; chunked or
fused cross-entropy, checkpointing, packing and context parallelism raise the cap.
\item Efficient reasoning optimises tokens per answer as much as cost per token.
\end{itemize}

% =====================================================================================
\section{Exercises}\label{sec:distill:exercises}

\begin{exercise}[Split balance]
Apply \ident{length_stats} to the training and validation partitions produced by
\ident{split_train_val}. Are they balanced in length? Propose a stratified split that guarantees
it, and explain why validation loss is length-sensitive.
\end{exercise}

\begin{exercise}[Without think tags]\label{ex:distill:nothink}
Repeat the distillation run with \code{use\_think=False} (the \texttt{--no\_think} flag of
\texttt{run.py}) and compare validation loss and, if possible, MATH-500 accuracy with the tagged
version. Why might the tags matter more for a model that must later switch between thinking and
non-thinking modes?
\end{exercise}

\begin{exercise}[Loss identity]
Prove that \ident{answer_only_loss} equals $-\frac{1}{n}\log p_\theta(y\mid x)$, and that the
token-level batched loss equals the average of per-example losses if and only if all answers have
the same length.
\end{exercise}

\begin{exercise}[Off-by-one]
Modify \ident{answer_only_loss} to start at \code{prompt\_len} and train the lab student. Measure
how often the trained student emits \texttt{<think>} as its first token.
\end{exercise}

\begin{exercise}[KD temperature]
Show that the gradient of $\tau^2\KL(p_T^{(\tau)}\|p_S^{(\tau)})$ with respect to a student logit
is $\tau(p_S^{(\tau)}-p_T^{(\tau)})$, and explain why the $\tau^2$ factor is needed when mixing it
with the hard-label cross-entropy.
\end{exercise}

\begin{exercise}[Top-$k$ storage]
For a 3{,}000-token trace and $V=151{,}936$, compute the storage for full bf16 distributions, for
top-$k$ with $k\in\{8,32,128\}$, and for text. Using the lab's \ident{sparse_kd_loss}, measure how
the loss approximation error depends on $k$ for a peaked and a flat teacher.
\end{exercise}

\begin{exercise}[Forward vs.\ reverse KL]
Fit a single Gaussian to a mixture of two Gaussians by minimising forward and reverse KL
numerically. Relate the result to \cref{fig:distill:kl} and to the behaviour of a small student
that cannot represent all of the teacher's reasoning strategies.
\end{exercise}

\begin{exercise}[Logit memory]
For Qwen3-0.6B at $T\in\{2\text{k},8\text{k},32\text{k}\}$, estimate total training memory
(weights, optimizer, activations with and without checkpointing, logits with and without
chunking). Which term dominates at each length?
\end{exercise}

\begin{exercise}[Packing for SFT]
Implement packed training for distillation using the block-diagonal mask of \cref{ch:rlsys}
(\ident{varlen_forward} is inference-only; write a differentiable variant). Verify that the packed
loss equals the padded loss with token-level normalisation.
\end{exercise}

\begin{exercise}[Cost per solved problem]
Given two students---A: 33\% accuracy, 2{,}100 tokens per answer; B: 30\% accuracy, 900 tokens per
answer---and a serving cost per token, which is cheaper per solved problem? At what accuracy would
B have to fall for A to win?
\end{exercise}

\begin{exercise}[Teacher generation throughput]
Using the napkin math of \cref{sec:distill:generate} and the serving model of \cref{ch:serving},
estimate how long one 8$\times$H100 node would need to generate the 36M teacher tokens with an
open-weight 32B dense teacher. What limits throughput?
\end{exercise}

% =====================================================================================
\section{Lab 30: Distilling a reasoning teacher}\label{sec:distill:lab}

\begin{labbox}{30}{Distilling a reasoning teacher}
\labdir{lab30_distill}
\labsection{Goals}
Build the complete hard-distillation pipeline of this chapter at toy scale on a CPU, then add
soft distillation and the systems pieces (batched losses, chunked cross-entropy, top-$k$ storage,
padding vs.\ packing). The same code drives a real Qwen3-0.6B run through \texttt{run.py}.

\labsection{Setup}
The ``teacher'' is a rule-based solver for multi-digit addition that writes one line per column
(\texttt{7+5+0=12}, carries included) inside \texttt{<think>}\ldots\texttt{</think>}, followed by
\texttt{\textbackslash boxed\{sum\}}---a miniature of R1-style traces with a verifiable answer. A character-level
tokenizer has atomic special tokens (\texttt{<user>}, \texttt{<assistant>}, \texttt{<think>},
\texttt{</think>}, \texttt{<eos>}, \texttt{<pad>}). The student is a tiny Qwen3-shaped model from
the core package.
\begin{lstlisting}[style=shell]
cd code
pytest -q labs/lab30_distill                    # grades starter.py
LAB_IMPL=solution pytest -q labs/lab30_distill  # reference
python labs/lab30_distill/bench.py --digits 2   # trace ablation (minutes on CPU)
\end{lstlisting}

\labsection{Tasks}
\begin{enumerate}
\item \textbf{Data formatting.} Implement \ident{format_distilled_answer} and
\ident{build_examples}: chat-wrapped prompt, raw target, end token, \code{prompt\_len};
skip malformed records.
\item \textbf{Filtering and splitting.} Implement \ident{length_stats} (mean, min/max with
indices, p50/p90/p99), \ident{filter_by_max_len} and a seeded \ident{split_train_val}.
\item \textbf{Losses.} Implement \ident{answer_only_loss} (tested against
$-\ident{sequence_logprob}/n$), \ident{collate_answer_only} and \ident{batch_answer_only_loss}
with both token- and sequence-level reductions, \ident{chunked_cross_entropy},
\ident{kd_loss} (forward/reverse, temperature), \ident{topk_teacher} and \ident{sparse_kd_loss}.
\item \textbf{Training.} Implement \ident{train_distillation}; the test requires the validation
loss to drop by 40\% within four epochs. A second test trains a small student to match a larger
teacher's distributions with \ident{kd_loss}.
\item \textbf{Systems helpers.} Implement \ident{padding_waste} (padded, length-sorted, packed
real-token fractions) and \ident{logit_storage_bytes} (dense vs.\ top-$k$).
\end{enumerate}

\labsection{Measurements}
\begin{itemize}
\item Run \texttt{bench.py} to train three students on the same problems: full trace with think
tags, trace without tags, and answer only. Report held-out exact-match accuracy and mean target
length (\cref{fig:distill:labresult}).
\item Measure the real-token fraction of padded vs.\ packed batches on a mix of 1-digit and 6-digit
problems, and predict the speedup packing would give.
\item Plot dense vs.\ top-$k$ storage for a 3{,}000-token trace as a function of $k$.
\end{itemize}

\labsection{Deliverables}
Passing tests; the ablation table with a short interpretation (does the student need the trace?
does it need the tags?); the storage plot.

\labsection{Stretch goals}
\begin{itemize}
\item Real scale: run \texttt{run.py} on Qwen3-0.6B-Base with the public DeepSeek-R1 MATH traces
(\code{--max\_len 2048 --epochs 3 --lr 1e-5}), then evaluate checkpoints on MATH-500 with your
Lab~21 harness.
\item Add rejection sampling: corrupt 20\% of the teacher's answers and compare training on all
traces vs.\ verified traces only.
\item Train with \code{batch\_size=8} using both reductions and compare learning curves; then
implement packed training (Exercise~30.9) and measure tokens/s.
\item Distil with a mixed loss $\alpha\,\text{CE}+(1-\alpha)\tau^2\text{KD}$ from a tiny trained
teacher model and compare sample efficiency with pure hard distillation.
\end{itemize}
\end{labbox}
